\documentclass[
    a4paper,
    cleardoubleempty,
    BCOR1cm
]{scrbook}

\input{header}

\usepackage{amssymb}
\usepackage{amsmath}
\usepackage{algorithm2e}
\usepackage{caption}
\usepackage{graphicx}
\usepackage{listings}
\usepackage{xcolor}

% this is adapted from from the source code of the ScholarCopilot paper published on arxiv.
\lstset{
    basicstyle=\ttfamily\small,      % 设定基础字体和字号
    columns=fullflexible,            % 更美观的宽度控制
    frame=single,                    % 单一边框
    breaklines=true,                 % 自动换行
    breakatwhitespace=true,          % 优先在空白处换行
    keepspaces=true,                 % 保留空格
    showspaces=false,                % 不显示空格标记
    showstringspaces=false,          % 字符串空格不显示特殊标记
    showtabs=false,                  % 不显示tab标记
    tabsize=4,                       % tab宽度
    rulecolor=\color{gray!50},       % 边框颜色
    backgroundcolor=\color{gray!5},  % 浅色背景
    language=[LaTeX]TeX,
    literate={\ \ }{{\ }}1,           % 改善空格显示
    numbers=none
}

\RestyleAlgo{ruled}
\LinesNumbered

\captionsetup[table]{position=bottom}

\DeclareMathOperator*{\argmax}{arg\,max}

\newcommand{\refsec}[1]{section \ref{#1}}
\newcommand{\reffig}[1]{\figurename{ \ref{#1}}}
\newcommand{\refeq}[1]{Eq. \ref{#1}}
\newcommand{\refalg}[1]{Algo. \ref{#1}}
\newcommand{\reftab}[1]{Table \ref{#1}}

\definecolor{commentgreen}{rgb}{0.5,0.5,0}
\newcommand{\note}[1]{\textcolor{commentgreen}{// #1}\\}% colored c-style comment
\newcommand{\mathnote}[1]{\textcolor{commentgreen}{\text{// #1}}}% for math env
\newcommand{\D}{\text{d}}% leibniz notation d for derivatives/integrals
\newcommand{\lp}[1]{\left({#1}\right)}% large parentheses, e.g. for vectors or fractions
\newcommand{\lpvar}[1]{\left\{{#1}\right\}}% large {} parentheses, e.g. for sets
\newcommand{\mono}[1]{\texttt{#1}}% monospace text
\newcommand{\ceil}[1]{\lceil{#1}\rceil}
\newcommand{\floor}[1]{\lfloor{#1}\rfloor}

% \title{LLM-Enhanced Scientific Literature Search with Active Disambiguation}
% \author{My Name \thanks{e-mail: my.name@uni-tuebingen.de}}
% \date{\today}
\begin{document}

\begin{tabular}{lr}
% 
\includegraphics[width=0.5\linewidth]{logo_sw} % logo bw
 
\includegraphics[width=0.5\linewidth]{UT_WBMW_Rot_4C} % logo red
 & \hspace{0.2\linewidth}
 \parbox{0.5\linewidth}{
   \large\bf\textsf{\color{rot}{Mathematisch-\\Naturwissenschaftliche\\Fakultät\\\\}}
   \hspace{-.144cm}\normalsize\textsf{\color{rot}{Lernbasierte Computer Vision}}
   \vspace{0.6cm}
 }
\end{tabular}

\vspace*{10ex}
Masterarbeit

{\huge\bf\textsf{LLM-Enhanced Scientific Literature Search with Active Disambiguation}}

\vspace*{30ex}

Eberhard Karls Universität Tübingen\\
Mathematisch-Naturwissenschaftliche Fakultät\\
Wilhelm-Schickard-Institut für Informatik\\
Lernbasierte Computer Vision\\
Niklas Munkes,~ \verb+niklas.munkes@student.uni-tuebingen.de+,~ 2025

\vspace*{5ex}

\begin{tabular}{@{}l@{\hspace{2em}}l}
  Bearbeitungszeitraum:& 14.04.-14.10.2025 \vspace*{5ex} \\
  Gutachter:& Prof. Dr. Andreas Geiger, Universität Tübingen\\
  Zweitgutachter:& Prof. Dr. Carsten Eickhoff, Universität Tübingen\\
  Betreuer:& Haoyu He, Universität Tübingen
\end{tabular}

\thispagestyle{empty}
\newpage

\chapter*{Selbstst\"andigkeitserkl\"arung}
Hiermit versichere ich, dass ich die vorliegende Masterarbeit selbst\"andig und
nur mit den angegebenen Hilfsmitteln angefertigt habe und dass alle Stellen,
die dem Wortlaut oder dem Sinne nach anderen Werken entnommen sind,
durch Angaben von Quellen als Entlehnung kenntlich gemacht worden sind.
Diese Masterarbeit wurde in gleicher oder \"ahnlicher Form in keinem anderen
Studiengang als Pr\"ufungsleistung vorgelegt.

\vspace*{8ex}
\hrule
\vspace*{2ex}
Niklas Munkes (Matrikelnummer 4269436), \today



\chapter*{Abstract}
TODO

\tableofcontents

\chapter{Introduction}\label{introduction}
TODO

% Introduction RAG
\section{Retrieval-Augmented Generation}\label{rag}
% traditional rag
Retrieval-Augmented Generation (RAG, \cite{lewis2021retrievalaugmentedgenerationknowledgeintensivenlp}) models are designed to generate text that is grounded in real-world factual knowlege. They have been widely adopted for knowledge-intensive NLP tasks \cite{shao2025reasonirtrainingretrieversreasoning} such as long-form question answering \cite{xia2025groundsentenceimprovingretrievalaugmented}, scientific question answering \cite{lála2023paperqaretrievalaugmentedgenerativeagent} and question synthesis \cite{xu2024kiwidatasetknowledgeintensivewriting}, generation with sentence-level citations \cite{zhang2024longciteenablingllmsgenerate, xia2025groundsentenceimprovingretrievalaugmented} and academic text generation \cite{asai2024openscholarsynthesizingscientificliterature, wang2025scholarcopilottraininglargelanguage}. RAG models usually consist of a generator, a pre-trained sequence-to-sequence (seq2seq) model (e.g. a Transformer-based \cite{vaswani2017attentionneed} decoder, see \refsec{related-work}), and a retriever with access to a retrieval corpus, which may be implemented as a dense vector index (e.g. a HNSW \cite{malkov2016efficientrobustapproximatenearest} index, see \refsec{hnsw}) \cite{lewis2021retrievalaugmentedgenerationknowledgeintensivenlp}. Thus they combine pre-trained paramentric memory with non-paramentric (retrieval-based) memory \cite{lewis2021retrievalaugmentedgenerationknowledgeintensivenlp}. 

The RAG architecture initially proposed by \cite{lewis2021retrievalaugmentedgenerationknowledgeintensivenlp}  consists of a query encoder $q$, a document index $d$ and a generator $p_\theta$ (see \reffig{fig:rag_architecture}) TODO

\begin{figure}[t]
    \includegraphics[width=\textwidth]{../images/rag_architecture.png}
    \caption{Overview of the RAG architecture \cite{lewis2021retrievalaugmentedgenerationknowledgeintensivenlp}. For finding the top-k documents $z_i$, the authors used Maximum Inner Product Search (MIPS).}\label{fig:rag_architecture}
\end{figure}

By using an external retrieval corpus, the knowledge RAG models use to inform their generation can directly be inspected and interpreted \cite{lewis2021retrievalaugmentedgenerationknowledgeintensivenlp}. Their TODO

% special challenges of scientific text (short)

% why bother with citations
citations build trust \cite{do2024facilitatinghumanllmcollaborationfactuality}
% hallucinations
% interpretability
% provenance
% ...
\chapter{Related Work}\label{related-work}
TODO % intro to this section

% Overview Transformer -> Llama -> Qwen
\section{Decoder Architectures}\label{decoders}
TODO

% Introduction RAG
\section{Retrieval-Augmented Generation}\label{rag}
% traditional rag
TODO

% special challenges of scientific text

% rag for scientific writing
Previous work on text generation with citations treated retrieval and generation as seperate stages during inference. For example, "Attribute First, then Generate" \cite{slobodkin2024attributefirstgeneratelocallyattributable} initially retrieves a set of relevant spans which are then grouped into coherent clusters which in turn are used to inform an iterative sentence-by-sentence generation step \cite{slobodkin2024attributefirstgeneratelocallyattributable}. Similarly, OpenScholar \cite{asai2024openscholarsynthesizingscientificliterature} first retrieves and reranks a set of papers relevant to the given query. These papers inform a subsequent generation step that produces a generated response $r_i$ \cite{asai2024openscholarsynthesizingscientificliterature}. This response is refined through iterative \emph{self-feedback} where during each iteration of the feedback loop $F=f_1,...,f_T$, the generator model generates sentences $f_j$ that describe potential improvements to the response and may also prompt the retrieval pipeline to retrieve additional papers, should $r_i$ require more information to generate an improved response $r_{i+1}$ \cite{asai2024openscholarsynthesizingscientificliterature}. Depending on the choice of $T$, OpenScholar still performs retrieval and generation largely seperately. ReClaim \cite{xia2025groundsentenceimprovingretrievalaugmented} TODO

ScholarCopilot \cite{wang2025scholarcopilottraininglargelanguage} is a RAG system developed for generating academic-style text with accurate citations. It is intended as a tool that assists the user with academic paper writing by providing iterative text generation with automated citation retrieval \cite{wang2025scholarcopilottraininglargelanguage}. ScholarCopilot also enables user interaction during generation \cite{wang2025scholarcopilottraininglargelanguage} by allowing the user to rewrite the generated text or forcefully trigger citation retrieval between the consecutive generation steps but this feature is beyond the scope of this thesis and is thus not utilized for the experiments. 
\chapter{Methods}\label{methods}
TODO % intro to this section

% dataset construction for retr + gen evaluation
\section{Dataset}
TODO \cite{He2024HDT}

% ScholarCopilot - architecture only!
\section{ScholarCopilot}\label{scholarcopilot}
Instead of treating retrieval and generation as separate processing stages (see \refsec{introduction}) ScholarCopilot \cite{wang2025scholarcopilottraininglargelanguage} integrates the information retrieval directly into the generation process. The model is trained to generate \emph{retrieval tokens} (\mono{[RET]}), that upon being generated, cause the model to interupt the generation process to retrieve a reference from the retrieval corpus \cite{wang2025scholarcopilottraininglargelanguage}. The ScholarCopilot retrieval corpus\footnote{\url{https://huggingface.co/datasets/TIGER-Lab/ScholarCopilot-Data-v1}} consists of abstracts of scientific papers, which are encoded by the ScholarCopilot model prior to training/inference and stored as a HNSW \cite{malkov2016efficientrobustapproximatenearest} index using Faiss \cite{douze2024faiss} for efficient retrieval (see \refsec{hnsw}). Standard ScholarCopilot uses the retrieved abstract to inform its subsequent generation by replacing the \mono{[RET]} tokens in the generated context with the corresponding abstracts. 

ScholarCopilot is intended to assist with academic writing by providing "[...] text completion and citation suggestions" (see the official demo\footnote{\url{https://huggingface.co/spaces/TIGER-Lab/ScholarCopilot}}\textsuperscript{,}\footnote{\url{https://www.youtube.com/watch?v=QlY7S52sWDA}} for an example). In the following we will thus assume that the model is used to generate a scientific paper with citations. Let
\begin{align}\label{eq:initial-context}
    Y^\mono{in}=y^\mono{in}_1,...,y^\mono{in}_l
\end{align}
denote the initial context given to the model. Since we seek to generate a scientific document, this context may consist of the title and abstract of the paper the model is supposed to generate. Let furthermore 
\begin{align}\label{eq:generated-tokens}
    Y^\mono{gen}_{1:i-1}=y^\mono{gen}_{1},...,y^\mono{gen}_{j-1},c_j,y^\mono{gen}_{j+1},...,y^\mono{gen}_{i-1}
\end{align}
denote a token sequence generated by the model where $c_j$ signifies that a \mono{[RET]} token was generated by the model at generation step $j$. Note that the model may generate the retrieval token at multiple steps, and that the number of generated retrieval tokens is equal to the number of citations in the generated paper. The representation $\hat{c}_j$ of a \mono{[RET]} token $c_j$ computed by the model serves as an aggregated representation of the token sequence generated up to the \mono{[RET]} token (e.g. $Y^\mono{gen}_{1:j-1}$ in \refeq{eq:generated-tokens}), in the same manner as BERT \cite{devlin2019bertpretrainingdeepbidirectional} (and also for example CF-ViT \cite{chen2022cfvitgeneralcoarsetofinemethod}) uses the \mono{[cls]} token as a representation of its input sequence. Thus $\hat{c}_j$ is unique for every $j$ since it is dependent of the corresponding left-hand context. 

When the ScholarCopilot generator $p_G$ generates a retrieval token (that is, $p_G(\ast|Y^\mono{in}Y^\mono{gen}_{1:i-1})$ is maximal for the \mono{[RET]} token, thus $\ast\leftarrow c_i$), the generation process is interrupted and instead, the representation $\hat{c}_i$ is passed to the retriever $p_R(z_i|\hat{c}_i)$ used to retrieve an encoded reference $z_i$ corresponding to an abstract
\begin{align}
    Y^\mono{ref}_i=y^\mono{ref}_{i,1},...,y^\mono{ref}_{i,k}
\end{align}
from the retrieval corpus. Before passing the newly generated context to the generator the retrieval token $c_i$ is replaced with the corresponding abstract to obtain an augmented generation output
\begin{align}
    \mathbb{Y}^\mono{gen}_{1:i}=y^\mono{gen}_{1},...,y^\mono{gen}_{i-1},y^\mono{ref}_{i,1},...,y^\mono{ref}_{i,k}
\end{align}
that is used to inform the sebsequent generation steps with the retrieved information. The generator $p_G(y_{i+1}|Y^\mono{in}\mathbb{Y}^\mono{gen}_{1:i})$ predicts the next token $y_{i+1}$ at generation step $i+1$ based on the given initial context $Y^\mono{in}$ and the augmented output $\mathbb{Y}^\mono{gen}_{1:i}$. See \refalg{alg:scholarcopilot} for a detailed overview of one generation step.

\begin{algorithm}[t]
    \captionsetup{labelfont={sc,bf}, labelsep=newline}
    \KwIn{The initial context $Y^\mono{in}$ and the previously generated sequence $Y^\mono{gen}_{1:i-1}$}
    \KwOut{The initial context $Y^\mono{in}$ and the new generated sequence $Y^\mono{gen}_{1:i}$}
    \TitleOfAlgo{ScholarCopilotGenerationStep($Y^\mono{in},Y^\mono{gen}_{1:i-1}$)}
    $y_i\leftarrow\argmax\limits_{y_i\in Y}p_G(y_i|Y^\mono{in}Y^\mono{gen}_{1:i-1})$\\
    \If{$y_i=\mono{[RET]}$}{
        $\hat{c}_i\leftarrow\mathcal{G}(y_i)$\\
        $z_i\leftarrow\argmax\limits_{z_i\in Z}p_R(z_i|\hat{c}_i)$\\
        $Y^\mono{ref}_i\leftarrow$ getAbstract$(z_i)$\\
        $\mathbb{Y}^\mono{gen}_{1:i}\leftarrow Y^\mono{gen}_{1:i-1}Y^\mono{ref}_i$\\
        \Return{$Y^\mono{in}$, $\mathbb{Y}^\mono{gen}_{1:i}$}
    }
    \Else{
        $Y^\mono{gen}_{1:i}\leftarrow Y^\mono{gen}_{1:i-1}y_i$\\
        \Return{$Y^\mono{in}$, $Y^\mono{gen}_{1:i}$}
    }
    \caption{Overview of one generation step of ScholarCopilot. Here $Y$ denotes the vocabulary (the set of all tokens), $Z$ the retrieval corpus (all encoded abstracts), $\mathcal{G}$ the ScholarCopilot backbone model (see \refsec{qwen}) and getAbstract$()$ a function that returns the unencoded abstract corresponding to $z_i$}\label{alg:scholarcopilot}
\end{algorithm}



% FAISS + HNSW
\subsection{HNSW}\label{hnsw}
TODO

\subsection{ScholarCopilot Training}\label{sc-training}
TODO

% Qwen 2.5 (as the reranker)
\section{Passage Reranking}\label{passage-reranking}
The goal of the \emph{Passage Retrieval} module (PR) introduced with this thesis is it to enhance the citation accuracy and generation quality of ScholarCopilot. As outlined in \refsec{scholarcopilot}, ScholarCopilot represents 

% Batched Self-Consistency
\subsection{Batched Self-Consistency}\label{batched-self-consistency}
TODO

% Min-Max Normalization of scores
\subsection{Min-Max Normalization}\label{normalization}
TODO


\chapter{Experiments}\label{experiments}
TODO

\section{Retrieval Evaluation}\label{retrieval-evaluation}
TODO
% parameters scale to recall k

\section{Generation Evaluation}\label{generation-evaluation}
TODO

% \section{Additional Experiments}\label{additional-experiments}
% TODO

\chapter{Limitations and Future Work}\label{limitations-future-work}
TODO

% train sc for direct passage retrieval
% build corpus from passages instead of abstracts
% passage reasoning
% 
\chapter{Conclusion}\label{conclusion}
In this thesis, I introduced the passage retrieval module that enables ScholarCopilot to cite fine-grained passages of scientific documents. While extensive experiments show that using ScholarCopilot with passage retrieval results in lower citation recall and does not increase the quality of the generation, the module is able to successfully retrieve the correct citation target, notably also in cases where ScholarCopilot fails. Qualitative analysis reveals the unreliable performance of only using an LLM for generation quality evaluation, suggesting that such evaluation cannot yet be conducted purely LLM-based with satisfactory results. Nevertheless, prompting the evaluating LLM to reason explicitly can consistently improve its evaluation performance, demonstrating potential for future research.


% \nocite{*} 
% \bibliographystyle{alpha}
\bibliography{bibliography}

\appendix
\chapter{Appendix}

\section{Hyperparameters}\label{hyperparameters}
TODO

\section{TODO}
TODO
% instructions
% numerical retr results
% PR with reasonir
% initial retrieval with reasonir instead of SC+passage_cutting

\end{document}

